% Part: normal-modal-logic
% Chapter: axioms-systems
% Section: proofs-modal-systems

\documentclass[../../../include/open-logic-section]{subfiles}

\begin{document}

\olfileid[hi]{nml}{prf}{prs}

\olsection{मोडल तंत्रों में प्रमाण}

अब हम $\Log{K}$ के अतिरिक्त अन्य मोडल तर्कशास्त्रीय तंत्रों के प्रमाणों पर
आते हैं।

\begin{prop}\ollabel{prop:S5facts}
निम्न प्रमाण्यता-परिणाम मान्य हैं:
  \begin{enumerate}
  \item $\Log{KT5} \Proves \Ax{B}$;
  \item $\Log{KT5} \Proves \Ax{4}$;
  \item $\Log{KDB4} \Proves \Ax{T}$;
  \item $\Log{KB4} \Proves \Ax{5}$;
  \item $\Log{KB5} \Proves \Ax{4}$;
  \item \ollabel{prop:S5facts-KT-D}$\Log{KT} \Proves \Ax{D}$.
  \end{enumerate}
\end{prop}

\begin{proof}
  हम प्रत्येक का प्रमाण प्रस्तुत करते हैं।
  \begin{enumerate}
    \item $\Log{KT5} \Proves \Ax{B}$:
      \begin{derivation}
        1. & $\Log{KT5} \Proves \Diamond!A \lif \Box\Diamond!A$ & \Ax{5}\\
        2. & $\Log{KT5} \Proves !A \lif \Diamond!A$ & $\Ax{T_\Diamond}$\\
        3. & $\Log{KT5} \Proves !A \lif \Box\Diamond!A$ & \PL, 2, 1.
      \end{derivation}
    \item $\Log{KT5} \Proves \Ax{4}$:
      \begin{derivation}
        1. &$\Log{KT5} \Proves \Diamond\Box!A \lif \Box\Diamond\Box!A$ & \Ax{5}
        जिसमें $p$ के स्थान पर $\Box!A$ है\\
        2. & $\Log{KT5} \Proves \Box!A \lif \Diamond\Box!A$ & $\Ax{T_\Diamond}$
        जिसमें $p$ के स्थान पर $\Box!A$ है\\
        3. & $\Log{KT5} \Proves \Box!A \lif \Box\Diamond\Box!A$ & \PL, 2, 1\\
        4. & $\Log{KT5} \Proves \Diamond\Box!A \lif \Box!A$ & $\Ax{5_\Diamond}$\\
        5. & $\Log{KT5} \Proves \Box\Diamond\Box!A \lif \Box\Box!A$ & \RK{}, 4 \\
        6. & $\Log{KT5} \Proves \Box!A \lif \Box\Box!A$ & \PL, 3, 5. \\
      \end{derivation}
    \item $\Log{KDB4} \Proves \Ax{T}$:
      \begin{derivation}
        1. & $\Log{KDB4} \Proves \Diamond\Box!A \lif !A $ & $\Ax{B_\Diamond}$ \\
        2. & $\Log{KDB4} \Proves \Box\Box!A \lif \Diamond\Box!A$ & $\Ax{D}$
        जिसमें $p$ के स्थान पर $\Box!A$ है\\
        3. & $\Log{KDB4} \Proves \Box\Box!A \lif !A$ & \PL, 1, 2\\
        4. & $\Log{KDB4} \Proves \Box!A \lif \Box\Box!A$ & \Ax{4} \\
        5. & $\Log{KDB4} \Proves \Box!A \lif !A$ & \PL, 4, 3. \\
      \end{derivation}
    \item $\Log{KB4} \Proves \Ax{5}$:
      \begin{derivation}
        1. & $\Log{KB4} \Proves \Diamond!A \lif \Box \Diamond\Diamond!A$ & \Ax{B}
        जिसमें $p$ के स्थान पर $\Diamond!A$ है\\
        2. & $\Log{KB4} \Proves \Diamond\Diamond !A \lif \Diamond !A$ &
        $\Ax{4_\Diamond}$ \\
        3. & $\Log{KB4} \Proves \Box\Diamond\Diamond !A \lif \Box \Diamond!A$ &
        \RK{}, 2 \\
        4. & $\Log{KB4} \Proves \Diamond!A \lif \Box\Diamond!A$ & \PL, 1, 3.
      \end{derivation}
    \item $\Log{KB5} \Proves \Ax{4}$:
      \begin{derivation}
        1. & $\Log{KB5} \Proves \Box!A \lif \Box\Diamond\Box !A$ & \Ax{B},
        जिसमें $p$ के स्थान पर $\Box!A$ है \\
        2. & $\Log{KB5} \Proves  \Diamond\Box!A \lif \Box !A$ &
        $\Ax{5_\Diamond}$\\
        3. & $\Log{KB5} \Proves  \Box\Diamond\Box!A \lif \Box\Box !A$ & \RK{}, 2 \\
        4. & $\Log{KB5} \Proves  \Box!A \lif \Box\Box!A$ & \PL, 1, 3.
      \end{derivation}
    \item $\Log{KT} \Proves \Ax{D}$:
      \begin{derivation}
        1. & $\Log{KT} \Proves \Box !A \lif !A$ & \Ax{T} \\
        2. & $\Log{KT} \Proves !A \lif \Diamond!A$ & $\Ax{T_\Diamond}$ \\
        3. & $\Log{KT} \Proves \Box !A \lif \Diamond !A$ &  \PL, 1, 2
      \end{derivation}
  \end{enumerate}
\end{proof}

\begin{defn}
  परंपरा के अनुसार हम \Log{S4} को तंत्र \Log{KT4} और \Log{S5} को तंत्र
  \Log{KTB4} परिभाषित करते हैं।
\end{defn}

निम्न प्रतिज्ञप्ति दिखाती है कि शास्त्रीय तंत्र \Log{S5} के अनेक तुल्य
अभिगृहीतीकरण हैं। यह आश्चर्यजनक नहीं, क्योंकि अभिगृहीतों के विभिन्न संयोजन
सभी तुल्यता संबंधों को अभिलक्षित करते हैं (देखें
\olref[frd][es5]{prop:equivalences})।

\begin{prop}\ollabel{prop:S5}
  $\Log{KTB4} = \Log{KT5} = \Log{KDB4} = \Log{KDB5}$.
\end{prop}

\begin{proof}
  अभ्यास।
\end{proof}

\begin{prob}
  \olref[nml][prf][prs]{prop:S5} सिद्ध कीजिए।
\end{prob}


\end{document}
